\section{Uniform Circular Motion}

One of the most important examples of motion is uniform circular motion.

A particle moving around a circle of radius

\[
R
\]

with angular velocity

\[
\omega
\]

has position

\[
x(t)=R\cos(\omega t),
\]

\[
y(t)=R\sin(\omega t).
\]

Differentiating,

\[
v_x
=
-R\omega\sin(\omega t),
\]

\[
v_y
=
R\omega\cos(\omega t).
\]

The speed is constant,

\[
v
=
R\omega,
\]

while the direction changes continuously.

The acceleration always points toward the centre of the circle,

\[
\boxed{
	a
	=
	\frac{v^2}{R}.
}
\]

Later, we shall discover that electrons in a magnetic field naturally execute
cyclotron motion, which is closely related to this simple example.

%%%%%%%%%%%%%%%%%%%%%%%%%%%%%%%%%%%%%%%%%%%%%%%%%%%%%%%%%%%%%%%%%%%%%%%%%%%%%%%%


\begin{bookfigure}{A parametric trajectory with instantaneous velocity tangent
to the curve and acceleration changing the velocity vector.}
\begin{tikzpicture}[scale=1.0]
  \draw[physics axis] (-0.3,0) -- (9.2,0) node[right] {$x$};
  \draw[physics axis] (0,-0.3) -- (0,4.7) node[above] {$y$};
  \draw[trajectory]
    plot[smooth] coordinates {(0.4,0.5) (1.5,1.25) (2.6,2.45)
      (3.9,2.9) (5.1,2.3) (6.3,1.4) (7.5,1.65) (8.5,3.0)};
  \fill[visualgold] (5.1,2.3) circle (2.2pt);
  \node[below,visual label] at (5.1,2.3) {$\mathbf r(t)$};
  \draw[physics vector] (5.1,2.3) -- (6.65,1.38)
    node[right,visual label] {$\mathbf v(t)$};
  \draw[force vector] (5.1,2.3) -- (4.55,3.65)
    node[above,visual label] {$\mathbf a(t)$};
\end{tikzpicture}
\end{bookfigure}

\begin{visualguidebox}
The position vector identifies the point on the curve. Velocity is tangent to
the trajectory. Acceleration need not point along the direction of motion.
\end{visualguidebox}
